\documentclass{article}

%Page format
\usepackage[table]{xcolor}
\usepackage{pdfpages}
\usepackage{fancyhdr}
\usepackage[margin=1in]{geometry}
\usepackage{tabularx}
%Math packages and custom commands
\usepackage{algpseudocode}
\usepackage{amsthm}
\usepackage{graphicx}
\usepackage{framed}
\usepackage[margin=1in]{geometry}
\usepackage{hyperref}
\usepackage{tikz}
\usepackage[utf8]{inputenc}
\definecolor{Gray}{gray}{0.9}
\usepackage{wrapfig}
\usepackage[margin=1in]{geometry}
\usepackage{mathtools,amsthm}
\usepackage{enumitem,amssymb}
\newtheoremstyle{case}{}{}{}{}{}{:}{ }{}
\theoremstyle{case}
\newtheorem{case}{Case}
\DeclareMathOperator{\R}{\mathbb{R}}
\DeclareMathOperator{\E}{\mathbb{E}}
\DeclareMathOperator{\Var}{\text{Var}}
\DeclareMathOperator{\Cov}{\text{Cov}}
\newcommand{\bvec}[1]{\mathbf{#1}}
\renewcommand{\P}{\mathbb{P}}
\newcommand{\norm}[2][2]{\| #2\|_{#1}}

\definecolor{shadecolor}{gray}{0.9}

\theoremstyle{definition}
\newtheorem*{answer}{Answer}
\newcommand{\note}[1]{\medskip \noindent{\textbf{NOTE:} #1}}
\newcommand{\hint}[1]{\medskip \noindent{\textit{HINT:} #1}}
\newcommand{\recall}[1]{\medskip{[\textbf{RECALL:} #1]}}
\newcommand{\motivation}[1]{\medskip \noindent{\textbf{MOTIVATION:} #1}}
\newcommand{\expect}[1]{\medskip \noindent{\fbox{\parbox{0.95 \textwidth}{\medskip \textbf{What we expect:} #1}}}}
\newcommand{\mysolution}[1]{\noindent{\begin{shaded}\textbf{Your Solution:}\ #1 \end{shaded}}}
\newlist{todolist}{itemize}{2}
\setlist[todolist]{label=$\square$}
\usepackage{pifont}
% use below to modify the size of equations
% usage is as follows: {fontsize}{math text size}{subscript size}{subsubscript size}
% if you wish to change the math font size, modify {math text size}{subscript size}{subsubscript size} for your chosen global font size {fontsize} which is declared at the beginning of the document
\DeclareMathSizes{10}{13}{13}{13}
\DeclareMathSizes{14}{17}{17}{17}
\DeclareMathSizes{12}{15}{15}{15}
\newcommand{\cmark}{\ding{51}}%
\newcommand{\xmark}{\ding{55}}%
\newcommand{\done}{\rlap{$\square$}{\raisebox{2pt}{\large\hspace{1pt}\cmark}}%
\hspace{-2.5pt}}
\newcommand{\wontfix}{\rlap{$\square$}{\large\hspace{1pt}\xmark}}
\graphicspath{ {./images/} }

\title{\textbf{CS221 Spring 2024: Artificial Intelligence:\\ Principles and Techniques} \\Homework 4: Controlling MountainCar}
\date{}

\chead{Route}
\rhead{\today}
\lfoot{}
\cfoot{CS221: Artificial Intelligence: Principles and Techniques --- Spring 2024}
\rfoot{\thepage}
\renewcommand{\headrulewidth}{0.4pt}
\renewcommand{\footrulewidth}{0.4pt}
\pagestyle{fancy}
\setlength{\parindent}{0pt}

\begin{document}

\maketitle

\begin{center}
\begin{tabular}{rl}
SUNet ID: & [your SUNet ID] \\
Name: & [your first and last name] \\
Collaborators: & [list all the people you worked with]
\end{tabular}
\end{center}
\newcolumntype{g}{>{\columncolor{Gray}}c}
\textit{By turning in this assignment, I agree by the Stanford honor code and declare
that all of this is my own work.} \\
% uncomment one of the below lines to make the text larger
\fontsize{12pt}{16pt}\selectfont
% \fontsize{14pt}{18pt}\selectfont
In route planning, the objective is to find the best way to get from point A to point B (think Google Maps).
  In this homework, we will build on top of the classic shortest path problem to allow for more powerful queries. For
  example, not only will you be able to explicitly ask for the shortest path from the Gates Building to the Green Library Coupa Cafe, but you can ask for the shortest path from Gates back to your dorm, stopping by the package center,
  gym, and the dining hall (in any order!) along the way. \\ 
  \\
\textbf{Before you get started, please read the Assignments section on the course website thoroughly}.



\section*{Problem 1: Value Iteration}





\begin{enumerate}



\item

  \mysolution{}
  


\item

  \mysolution{}

\end{enumerate}

\newpage



\end{document}